El estudio de factibilidad constituye una evaluación integral orientada a determinar si un proyecto propuesto —en este caso, el desarrollo e implementación del sistema de administración y gestión de comunidades residenciales \textbf{DOMU}— resulta técnica, operativa y económicamente viable. El objetivo principal de este análisis es estimar las posibilidades de éxito del proyecto antes de su ejecución, minimizando así los riesgos asociados a la inversión de tiempo, recursos y esfuerzos.

A través de esta evaluación, se busca proporcionar una base sólida para la toma de decisiones estratégicas, garantizando que \textbf{DOMU} no solo sea viable desde el punto de vista conceptual, sino también aplicable y rentable en contextos reales. En esta sección, se abordarán tres dimensiones fundamentales de la viabilidad del proyecto: la \textbf{factibilidad técnica}, la \textbf{factibilidad operativa} y, especialmente, la \textbf{factibilidad económica}, cuyo análisis permitirá establecer si los beneficios esperados justifican los costos asociados al desarrollo e implementación del sistema.

\section{Factibilidad Técnica}
Esta sección evalúa la viabilidad técnica del proyecto considerando tanto los recursos de desarrollo como la infraestructura necesaria para la operación en producción.

\begin{longtable}{|p{0.38\textwidth}|p{0.18\textwidth}|p{0.36\textwidth}|}
    \caption{Especificaciones del software requerido para el desarrollo.}
    \label{tab:software_desarrollo} \\
    \hline
    \multicolumn{3}{|c|}{\textbf{Especificación del software requerido en desarrollo}} \\
    \hline
    \textbf{Software} & \textbf{Versión Mínima} & \textbf{Propósito en el Proyecto} \\
    \hline
    \endhead
    \hline
    \multicolumn{3}{|r|}{\small\textit{Continúa en la siguiente página...}} \\
    \hline
    \endfoot
    \hline
    \endlastfoot

    \multicolumn{3}{|l|}{\textit{Entorno y Control de Versiones}} \\ \hline
    Sistema operativo PC & macOS, Windows 11 & Entorno compatible de desarrollo \\
    Git & 2.50+ & Sistema de control de versiones \\
    GitHub & Web & Plataforma de repositorios Git \\ \hline

    \multicolumn{3}{|l|}{\textit{IDEs y Editores de Código}} \\ \hline
    Visual Studio Code & 1.101+ & Editor de código (enfocado a front-end) \\
    IntelliJ IDEA & 2024.1+ & IDE de Java (enfocado a back-end) \\ \hline

    \multicolumn{3}{|l|}{\textit{Tecnologías Back-End}} \\ \hline
    Java JDK & 21 & Kit de desarrollo de Java \\
    Spring Boot & 3.x & Framework de aplicaciones Java \\ \hline

    \multicolumn{3}{|l|}{\textit{Tecnologías Front-End}} \\ \hline
    Node.js + NVM & LTS & Entorno de ejecución y gestión de versiones \\
    React / React Native / Expo & Actual & Frameworks para interfaces de usuario \\ \hline

    \multicolumn{3}{|l|}{\textit{Gestión de Datos y API}} \\ \hline
    MySQL & 8.0+ & Motor de base de datos relacional \\
    DBeaver & 25.1+ & Interfaz gráfica para MySQL \\
    Postman & v11+ & Cliente para pruebas de API REST \\ \hline

    \multicolumn{3}{|l|}{\textit{Contenerización}} \\ \hline
    Docker & 20.10.14+ & Plataforma de contenerización \\ \hline

    \multicolumn{3}{|l|}{\textit{Diseño y Navegación}} \\ \hline
    Figma & Web & Herramienta de diseño y prototipado \\
    Google Chrome & Última & Navegador web para depuración \\ \hline
\end{longtable}

\begin{table}[H]
    \centering
    \caption{Especificaciones del hardware requerido para el desarrollo.}
    \label{tab:hardware_desarrollo}
    \begin{tabular}{|p{0.42\textwidth}|p{0.50\textwidth}|}
    \hline
    \multicolumn{2}{|c|}{\textbf{Especificación hardware requerido en desarrollo}}\\
    \hline
    \textbf{Especificación hardware} & \textbf{Detalle mínimo} \\
    \hline
    Dispositivos para desarrollo & MacBook Pro M1 (8-core), 16~GB RAM, SSD~500~GB; HP Victus 16 (i7-10750H), 16~GB RAM, SSD~1~TB \\
    \hline
    Pistola escáner de QR & Capaz de leer códigos QR (ISO/IEC 18004), conexión USB o Bluetooth. \\
    \hline
    \end{tabular}
\end{table}

\begin{table}[H]
    \centering
    \caption{Especificaciones del hardware para los servidores de desarrollo y producción.}
    \label{tab:hardware_servidores}
    \begin{tabular}{|p{0.42\textwidth}|p{0.50\textwidth}|}
    \hline
    \multicolumn{2}{|c|}{\textbf{Especificación hardware servidor de desarrollo y producción}}\\
    \hline
    \textbf{Especificación} & \textbf{Detalle} \\
    \hline
    Servidor de desarrollo & Intel Core i5, 16~GB RAM, SSD; Windows~11. Pruebas locales y Docker. \\
    \hline
    Servidor de producción & Xeon multi-core 2.0~GHz, 32~GB RAM, SSD~NVMe 1~TB. \\
    \hline
    \end{tabular}
\end{table}

% --- NUEVA SECCIÓN AGREGADA PARA CUMPLIR CON LA PROFESORA ---
\subsection{Viabilidad de la Infraestructura en Producción y Mantenimiento}

Para dar cumplimiento a los requerimientos de operación continua una vez que el sistema sea entregado, se ha definido una estrategia técnica que asegura la disponibilidad y mantenibilidad del software en un entorno productivo:

\begin{itemize}
    \item \textbf{Disponibilidad del Servidor (SLA):} La infraestructura de producción se alojará en un entorno controlado (VPS o Cloud) que garantice un Acuerdo de Nivel de Servicio (SLA) del \textbf{99,0\% mensual}, alineado con el requerimiento no funcional RNF\_01. Esto se logrará mediante servidores con redundancia eléctrica y de red.
    
    \item \textbf{Estrategia de Mantenimiento y Despliegue:} El uso de \textbf{Docker} permite encapsular la aplicación y sus dependencias, facilitando la corrección de errores (hotfixes) y el despliegue de nuevas versiones sin generar conflictos con el sistema operativo base. Esto asegura que el mantenimiento correctivo y evolutivo sea ágil y reproducible.
    
    \item \textbf{Respaldo y Recuperación:} Para mitigar riesgos de pérdida de información, se configuran rutinas automáticas (\textit{cron jobs}) que ejecutan un respaldo diario (\textit{backup}) de la base de datos MySQL, garantizando la integridad de la información financiera y de accesos de la comunidad.
\end{itemize}

\textbf{Conclusión de la factibilidad técnica:} Con base en las especificaciones de software, hardware de desarrollo, la estrategia de servidores en producción y los planes de mantenimiento, se concluye que el proyecto \textbf{DOMU} es técnicamente factible.

El equipo de desarrollo cuenta con los recursos tecnológicos necesarios y un ecosistema de software moderno (Java, Spring Boot, React, MySQL y Docker). En cuanto al hardware, los componentes físicos requeridos para la operación (pistola escáner) son de fácil adquisición y alta compatibilidad.

Considerando la madurez de las tecnologías y la planificación de la infraestructura productiva, se afirma que el sistema es completamente viable desde el punto de vista técnico.

\section{Factibilidad Operativa}

La factibilidad operativa del sistema \textbf{DOMU} evalúa si los usuarios finales —administradores de comunidades, conserjes, comités, residentes y personal de servicio— podrán adoptar y utilizar eficazmente la plataforma en su entorno cotidiano. También considera la disposición de los actores involucrados para participar activamente y las competencias mínimas necesarias para su uso.

\subsection{Disponibilidad tecnológica y competencias de los usuarios}

El éxito operativo de una plataforma como \textbf{DOMU} está estrechamente vinculado al nivel de acceso a la tecnología y a las competencias digitales de sus usuarios. En este contexto, las condiciones actuales de conectividad en las regiones del Biobío y La Araucanía resultan favorables. De acuerdo con datos de Subtel e INE, más del 80~\% de los hogares urbanos en estas regiones cuenta con conexión a Internet fija o móvil. Esto permite establecer que tanto residentes como miembros del comité poseen el equipamiento necesario para interactuar con la plataforma.

En lo que respecta a los administradores y conserjes, estos ya utilizan herramientas digitales básicas (WhatsApp, correo). La principal herramienta física para el registro en conserjería será la \textbf{pistola de escáner QR}, un dispositivo de uso simple que facilita el registro rápido y sin errores.

Por su parte, el personal de aseo y mantención solo requiere realizar acciones simples dentro del sistema, funciones que no exigen habilidades técnicas avanzadas y que pueden ser fácilmente abordadas mediante capacitaciones breves.

\subsection{Usabilidad del sistema y apoyo a la adopción}

Uno de los factores críticos para la adopción exitosa es la facilidad de uso. En el caso de \textbf{DOMU}, la interfaz será diseñada priorizando la simplicidad y claridad visual. Además, el sistema contempla compatibilidad multiplataforma (móvil y web). Esta accesibilidad técnica se complementa con elementos de apoyo integrados, como guías interactivas y ayuda contextual.

\subsection{Aceptación y motivación}

La disposición de los usuarios depende del valor percibido. \textbf{DOMU} responde a problemáticas concretas: desorganización en visitas, falta de trazabilidad en encomiendas y gestión financiera opaca. Al ofrecer soluciones automatizadas, el sistema entrega beneficios tangibles: los administradores mejoran el control, los conserjes reducen su carga manual y los residentes obtienen transparencia y comodidad.

\subsection*{Conclusión de la factibilidad operativa}

A partir del análisis realizado, se concluye que el sistema \textbf{DOMU} es operativamente factible. Las zonas de implementación presentan condiciones tecnológicas adecuadas y los usuarios poseen las competencias digitales mínimas. La estructura del sistema, diseñada con criterios de usabilidad, reduce las barreras de entrada y promueve un uso eficiente.

\section{Factibilidad Económica}

La factibilidad económica determina si el desarrollo e implementación resultan financieramente viables. Esta sección identifica y cuantifica los costos asociados y proyecta los beneficios para fundamentar la sostenibilidad del proyecto.

En esta evaluación se considerarán los costos directos (RRHH, licencias, hosting, hardware), así como la depreciación de los activos utilizados (equipos de desarrollo), calculada proporcionalmente al periodo de desarrollo (4 meses).

\begin{table}[H]
    \centering
    \caption{Costos de recursos técnicos de desarrollo (Software y Hardware).}
    \label{tab:costos_recursos_tecnicos}
    {\small\setlength{\tabcolsep}{4pt}
    \begin{tabular}{|p{0.23\textwidth}|c|r|c|r|c|r|}
    \hline
    \textbf{Recurso} & \textbf{Cant.} & \textbf{Valor (CLP)} & \textbf{Depr. (mes)} & \textbf{Valor mensual} & \textbf{Meses} & \textbf{Total (CLP)}\\
    \hline
    \multicolumn{7}{|l|}{\textbf{Software}}\\ \hline
    Visual Studio Code & 2 & \$0 & & & 4 & \$0\\
    Docker & 2 & \$0 & & & 4 & \$0\\
    Git & 2 & \$0 & & & 4 & \$0\\
    GitHub & 2 & \$0 & & & 4 & \$0\\
    Sistema op. Windows & 1 & \$229.999 & 36 & \$6.389 & 4 & \$25.555\\
    Sistema op. macOS & 1 & \$0 & & & 4 & \$0\\
    Java JDK & 2 & \$0 & & & 4 & \$0\\
    Javalin & 2 & \$0 & & & 4 & \$0\\
    Node.js + NVM & 2 & \$0 & & & 4 & \$0\\
    React / Expo & 2 & \$0 & & & 4 & \$0\\
    MySQL & 1 & \$0 & & & 4 & \$0\\
    Postman & 2 & \$0 & & & 4 & \$0\\
    Google Chrome & 2 & \$0 & & & 4 & \$0\\
    DBeaver & 2 & \$0 & & & 4 & \$0\\
    Figma & 1 & \$0 & & & 4 & \$0\\ \hline
    \multicolumn{6}{|r|}{\textbf{Total software}} & \$25.555\\
    \hline \hline
    \multicolumn{7}{|l|}{\textbf{Hardware}}\\ \hline
    MacBook Pro M1 & 1 & \$1.200.000 & 36 & \$33.333 & 4 & \$133.332\\
    HP Victus 16 & 1 & \$709.990 & 36 & \$19.722 & 4 & \$78.888\\
    Pistola escáner QR & 1 & \$22.525 & 24 & \$939 & 4 & \$3.756\\ \hline
    \multicolumn{6}{|r|}{\textbf{Total hardware}} & \$215.976\\
    \hline \hline
    \multicolumn{6}{|r|}{\textbf{Total recursos técnicos}} & \$241.531\\
    \hline
    \end{tabular}
    }
\end{table}

A continuación, la Tabla~\ref{tab:costos_servicios} detalla los costos de adquisición y mantenimiento de dominio web y hosting, fundamentales para la operación continua.

\begin{table}[H]
    \centering
    \caption{Costos de Hosting y Dominio.}
    \label{tab:costos_servicios}
    \begin{tabular}{|l|r|r|}
    \hline
    \textbf{Ítem} & \textbf{Mensual (CLP)} & \textbf{Anual (CLP)}\\
    \hline
    Dominio & \$832 & \$9.990\\
    Hosting & \$13.500 & \$162.000\\
    \hline
    \textbf{Total} & \textbf{\$14.332} & \textbf{\$171.990}\\
    \hline
    \end{tabular}
\end{table}

En la Tabla~\ref{tab:costos_rrhh} se presenta la estimación del costo de recursos humanos (costo de oportunidad) basado en tarifas de mercado para un periodo de 4 meses con dedicación parcial.

\begin{table}[H]
    \centering
    \caption{Cálculo costo de recursos humanos de desarrollo.}
    \label{tab:costos_rrhh}
    {\small\setlength{\tabcolsep}{3pt}
    \renewcommand{\arraystretch}{1.1}
    \begin{tabular}{|p{0.34\textwidth}|c|c|c|r|r|}
    \hline
    \textbf{Cargo} & \textbf{Cant.} & \textbf{Meses} & \shortstack{\textbf{Horas}\\\textbf{estim.}} & \textbf{Costo x hora} & \textbf{Total (CLP)}\\
    \hline
    Desarrollador Frontend & 1 & 4 & 160 & \$18.250 & \$2.920.000\\
    Desarrollador Backend  & 1 & 4 & 160 & \$20.000 & \$3.200.000\\
    Diseñador UX/UI        & 1 & 1 &  40 & \$15.000 &   \$600.000\\
    Tester / QA            & 1 & 1 &  40 & \$15.000 &   \$600.000\\
    Gestor de Proyecto     & 1 & 4 & 160 & \$21.666 & \$3.466.560\\
    \hline
    \multicolumn{5}{|r|}{\textbf{Total}} & \textbf{\$10.786.560}\\
    \hline
    \end{tabular}
    }
\end{table}

La Tabla~\ref{tab:costos_suministros} estima los gastos de servicios básicos (Internet y electricidad) necesarios para el desarrollo remoto.

\begin{table}[H]
    \centering
    \caption{Cálculo del costo de suministros de desarrollo.}
    \label{tab:costos_suministros}
    \begin{tabular}{|l|c|r|c|r|}
    \hline
    \textbf{Material} & \textbf{Cant.} & \textbf{Costo mensual} & \textbf{Meses} & \textbf{Total (CLP)}\\
    \hline
    Plan Internet & 2 & \$16.990 & 4 & \$67.960\\
    Boletas de electricidad & 2 & \$12.000 & 4 & \$48.000\\
    \hline
    \textbf{Total} & & \textbf{\$28.990} & & \textbf{\$115.960}\\
    \hline
    \end{tabular}
\end{table}

Finalmente, la Tabla~\ref{tab:costos_resumen} resume la inversión inicial total (año cero).

\begin{table}[H]
    \centering
    \caption{Resumen de Costos de Desarrollo.}
    \label{tab:costos_resumen}
    \begin{tabular}{|l|r|}
    \hline
    \multicolumn{2}{|c|}{\textbf{Resumen de Costos de Desarrollo}}\\
    \hline
    Duración & 4 meses\\
    \hline
    Recursos Técnicos & \$241.531\\
    \hline
    Recursos Humanos & \$10.786.560\\
    \hline
    Suministros & \$115.960\\
    \hline
    \textbf{Total} & \textbf{\$11.144.051}\\
    \hline
    \end{tabular}
\end{table}

\subsection{Flujo de Caja}

Se estima un valor de suscripción de \$90.000 CLP mensual por comunidad. La proyección de ingresos considera un crecimiento escalonado hasta alcanzar 25 comunidades el primer año, apoyado por marketing digital (Meta Ads, Google Ads).

\begin{table}[H]
    \centering
    \caption{Estimación de ingresos mensuales (primer año).}
    \label{tab:ingresos_y1}
    \begin{tabular}{|c|c|r|r|}
    \hline
    \textbf{Mes} & \textbf{Comunidades} & \textbf{Valor suscripción} & \textbf{Ingreso mensual (CLP)}\\
    \hline
    1 & 3 & \$90.000 & \$270.000\\
    2 & 5 & \$90.000 & \$450.000\\
    3 & 8 & \$90.000 & \$720.000\\
    4 & 8 & \$90.000 & \$720.000\\
    5 & 10 & \$90.000 & \$900.000\\
    6 & 11 & \$90.000 & \$990.000\\
    7 & 13 & \$90.000 & \$1.170.000\\
    8 & 15 & \$90.000 & \$1.350.000\\
    9 & 19 & \$90.000 & \$1.710.000\\
    10 & 21 & \$90.000 & \$1.890.000\\
    11 & 22 & \$90.000 & \$1.980.000\\
    12 & 25 & \$90.000 & \$2.250.000\\
    \hline
    \textbf{Total} & & & \textbf{\$14.400.000}\\
    \hline
    \end{tabular}
\end{table}

Para los costos operativos, se considera publicidad, soporte (practicantes TI) y mantenimiento técnico.

\begin{table}[H]
    \centering
    \caption{Estimación de Costos Operativos del proyecto.}
    \label{tab:costos_operativos_y1}
    \begin{tabular}{|l|r|}
    \hline
    \multicolumn{2}{|c|}{\textbf{Estimación de Costos Operativos}}\\
    \hline
    \textbf{Ítem} & \textbf{Valor Anual (12 meses)}\\
    \hline
    Publicidad (Meta y Google Ads) & \$3.600.000\\
    \hline
    Soporte y Atención & \$6.348.000\\
    \hline
    Mantenimiento Técnico & \$8.400.000\\
    \hline
    Dominio y Hosting & \$171.990\\
    \hline
    \textbf{Total Costos Operativos} & \textbf{\$18.519.990}\\
    \hline
    \end{tabular}
\end{table}

La proyección de suscripciones para los años siguientes se basa en una tasa compuesta de crecimiento anual (CAGR) del 50\%.

$$
\text{Suscripciones en año } n = S_0 \times (1 + r)^n
$$
Donde $S_0 = 25$ y $r = 0.5$.

\begin{table}[H]
    \centering
    \caption{Proyección de Suscripciones por Año (CAGR 50\%).}
    \label{tab:proyeccion_suscripciones}
    \begin{tabular}{|c|c|c|}
    \hline
    \multicolumn{3}{|c|}{\textbf{Proyección de Suscripciones por Año}}\\
    \hline
    \textbf{Año} & \textbf{Cálculo} & \textbf{Suscripciones Proyectadas}\\
    \hline
    1 & 25 & 25\\
    \hline
    2 & $25 \times (1 + 0.50)^1$ & 38\\
    \hline
    3 & $25 \times (1 + 0.50)^2$ & 56\\
    \hline
    4 & $25 \times (1 + 0.50)^3$ & 84\\
    \hline
    5 & $25 \times (1 + 0.50)^4$ & 127\\
    \hline
    \end{tabular}
\end{table}

A continuación, se presenta el flujo de caja proyectado a 5 años.

\begin{table}[H]
    \centering
    \caption{Flujo de caja del proyecto (proyectado a 5 años).}
    \label{tab:flujo_de_caja}
    \footnotesize % Reducimos un poco más el tamaño para que quepa bien
    \setlength{\tabcolsep}{3pt} % Reduce el espacio entre columnas
    \begin{tabularx}{\textwidth}{|X|r|r|r|r|r|r|}
    \hline
     & \textbf{Año 0} & \textbf{Año 1} & \textbf{Año 2} & \textbf{Año 3} & \textbf{Año 4} & \textbf{Año 5} \\
    \hline
    Suscripciones estimadas & 0 & 25 & 38 & 56 & 84 & 127 \\
    \hline
    Ingresos anuales (CLP) & \$0 & \$14.4M & \$41.0M & \$60.5M & \$90.7M & \$137.2M \\
    \hline
    Costos Operativos (CLP) & \$0 & \$18.5M & \$18.5M & \$18.5M & \$18.5M & \$18.5M \\
    \hline
    Inversión Inicial (CLP) & \$11.1M & \$0 & 0 & 0 & 0 & 0 \\
    \hline
    \textbf{Total (Flujo neto)} & \textbf{-\$11.1M} & \textbf{-\$4.1M} & \textbf{\$22.5M} & \textbf{\$42.0M} & \textbf{\$72.2M} & \textbf{\$118.6M} \\
    \hline
    \end{tabularx}
\end{table}

\subsection{Cálculo del V.A.N.}

$$
VAN = \sum_{t=1}^{n} \frac{F_t}{(1 + k)^t} - I_0
$$

\begin{table}[H]
    \centering
    \caption{Términos de la fórmula de VAN.}
    \label{tab:van_terminos}
    \begin{tabular}{|c|l|}
    \hline
    \textbf{Término} & \textbf{Significado} \\
    \hline
    $t$ & Intervalo de tiempo \\
    \hline
    $n$ & Duración en años \\
    \hline
    $I_0$ & Inversión inicial (t=0) \\
    \hline
    $k$ & Tasa de descuento \\
    \hline
    $F_t$ & Flujos de caja obtenidos en el intervalo de tiempo $t$ \\
    \hline
    \end{tabular}
\end{table}

A continuación, se calculará el VAN con una tasa de descuento del 10\%.

\begin{table}[H]
    \centering
    \caption{Cálculo del VAN.}
    \label{tab:van_calculo}
    \setlength{\tabcolsep}{4pt}
    \begin{tabular}{|l|l|r|}
    \hline
    \textbf{Año} & \textbf{Cálculo (k = 10\%)} & \textbf{Flujo de caja (PV)} \\
    \hline
    Año 0 & $-\$11.144.051 / (1+0.10)^{0}$ & \textbf{-\$11.144.051} \\
    \hline
    Año 1 & $-\$4.119.990 / (1+0.10)^{1}$  & \textbf{-\$3.745.445} \\
    \hline
    Año 2 & $\$22.520.010 / (1+0.10)^{2}$  & \textbf{\$18.611.579} \\
    \hline
    Año 3 & $\$41.960.010 / (1+0.10)^{3}$  & \textbf{\$31.525.166} \\
    \hline
    Año 4 & $\$72.200.010 / (1+0.10)^{4}$  & \textbf{\$49.310.298} \\
    \hline
    Año 5 & $\$118.640.010 / (1+0.10)^{5}$ & \textbf{\$73.662.292} \\
    \hline
    \end{tabular}
\end{table}

\textbf{VAN (10\%)} = -11.144.051 + (-3.745.445) + 18.611.579 + 31.525.166 + 49.310.298 + 73.662.292 = \textbf{\$158.219.839}

En este caso, el VAN obtenido es positivo (\textbf{\$158.219.839}), lo que indica que el proyecto es viable desde una perspectiva económica.

\section{Conclusión de Factibilidad}
Gracias al análisis realizado, se concluye que el proyecto es económicamente viable. El VAN obtenido asciende a \textbf{\$158.219.839} utilizando una tasa de descuento del 10\%. Este resultado indica que el proyecto no solo recupera su inversión inicial, sino que genera un excedente significativo. Asimismo, el flujo de caja proyectado alcanza resultados positivos a partir del segundo año, reforzando la sostenibilidad financiera a largo plazo.